%------------------------------------------------------------------------------%
\documentclass[fleqn,aspectratio=169,12pt]{beamer}

\usepackage{gaal}

\title{Representação Matricial de Equações}

%------------------------------------------------------------------------------%
\begin{document}

\addtitle

%------------------------------------------------------------------------------%
\section{Representação Matricial de Equações}

%------------------------------------------------------------------------------%
\begin{frame}{Representação Matricial de Equações}
  É conveniente escrever equações utilizando matrizes
\end{frame}

%------------------------------------------------------------------------------%
\begin{frame}{Equação da Reta no Plano}
  \onslide<+->{Definindo}
  \[
    \onslide<+->{v = \begin{bmatrix} x \\ y \end{bmatrix}}
    \qquad
    \onslide<+->{B = \begin{bmatrix} a \\ b \end{bmatrix} \neq \vec{0}}
    \qquad
    \onslide<+->{d \in \R}
  \]

  \onslide<+->{Podemos escrever a equação da reta no plano}
  \begin{align*}
    \onslide<+->{B^t v & = d} \\
    \onslide<+->{\begin{bmatrix} a & b \end{bmatrix} \begin{bmatrix} x \\ y \end{bmatrix} & = d}\\
    \onslide<+->{ax + by              & = d}
  \end{align*}
\end{frame}

%------------------------------------------------------------------------------%
\begin{frame}{Planos em Forma Matricial}
  \onslide<+->{Definindo}
  \[
    \onslide<+->{v = \begin{bmatrix} x \\ y \\ z \end{bmatrix}}
    \qquad
    \onslide<+->{B = \begin{bmatrix} a \\ b \\ c \end{bmatrix} \neq \vec{0}}
    \qquad
    \onslide<+->{d \in \R}
  \]

  \onslide<+->{Podemos escrever a equação do plano no espaço}
  \begin{align*}
    \onslide<+->{B^tv & = d} \\
    \onslide<+->{ax + by +cz          & = d}
  \end{align*}
\end{frame}

%------------------------------------------------------------------------------%
\begin{frame}{Formas Quadráticas}
  \onslide<+->{Definindo}
  \[
    \onslide<+->{v = \begin{bmatrix} x \\ y \end{bmatrix}}
    \qquad
    \onslide<+->{A = \begin{bmatrix}
      a & b \\
      b & c
    \end{bmatrix} }
  \]

  \onslide<+->{Podemos escrever a expressão do segundo grau}
  \begin{align*}
    \onslide<+->{v^tAv}
    \onslide<+->{& = \begin{bmatrix}
      x & y
    \end{bmatrix}
    \begin{bmatrix}
      a & b \\
      b & c
    \end{bmatrix}
    \begin{bmatrix}
      x \\
      y
    \end{bmatrix}}
    \onslide<+->{\;=\;
    \begin{bmatrix}
      x & y
    \end{bmatrix}
    \begin{bmatrix}
      ax + by \\
      bx + cy
    \end{bmatrix}}
    \\
    \onslide<+->{& = x(ax + by) + y(bx + cy)}
    \\
    \onslide<+->{& = ax^2 + 2bxy + cy^2}
  \end{align*}
\end{frame}

%------------------------------------------------------------------------------%
\begin{frame}{Equação Geral de Segundo Grau}
  \onslide<+->{Definindo}
  \[
    \onslide<+->{v = \begin{bmatrix} x \\ y \end{bmatrix}}
    \qquad
    \onslide<+->{A = \begin{bmatrix}
      a & b \\
      b & c
    \end{bmatrix}}
    \qquad
    \onslide<+->{B = \begin{bmatrix} d \\ e \end{bmatrix}}
    \qquad
    \onslide<+->{f \in \R}
  \]

  \onslide<+->{Qualquer equação de primeiro ou segundo grau pode ser escrita como}
  \[
    \onslide<+->{v^tAv \;+\; B^tv \;+\; f = 0}
  \]
  \onslide<+->{Forma expandida
  \[
    ax^2 + 2bxy + cy^2 + dx + ey + f = 0
  \]}
\end{frame}

%------------------------------------------------------------------------------%
%------------------------------------------------------------------------------%
\begin{question}
  Expanda a equação
  \(
    v^tAv + B^tv + f = 0
  \)
  com
  \[
    v = \begin{bmatrix} x \\ y \end{bmatrix}
    \qquad
    A = \begin{bmatrix}
      2 & -3 \\
      -3 & 0
    \end{bmatrix}
    \qquad
    B = \begin{bmatrix} 4 \\ -7 \end{bmatrix}
    \qquad
    f = -1
  \]
\end{question}

%------------------------------------------------------------------------------%
\begin{resolution}{Solução}
  \begin{align*}
    \onslide<+->{v^tAv + B^tv + f & = 0 }
    \\[2mm]
    \onslide<+->{\begin{bmatrix} x & y \end{bmatrix}
    \begin{bmatrix}
      2 & -3 \\
      -3 & 0
    \end{bmatrix}
    \begin{bmatrix} x \\ y \end{bmatrix}
    +
    \begin{bmatrix} 4 &  -7 \end{bmatrix}
    \begin{bmatrix} x \\ y \end{bmatrix}
    + (-1)
    & = 0}
    \\[2mm]
    \onslide<+->{2x^2 + 2(-3) xy + 0y^2}
    \onslide<+->{+ 4x + (-7)y - 1}
    \onslide<+->{& = 0}
    \\[2mm]
    \onslide<+->{2x^2 -6xy + 4x -7y & = 1}
  \end{align*}
\end{resolution}

%------------------------------------------------------------------------------%
\endinput

%------------------------------------------------------------------------------%
\begin{question}
  Expanda a equação
  \(
    v^tAv + B^tv + f = 0
  \)
  com
  \[
    v = \begin{bmatrix} x \\ y \\ z \\ w \end{bmatrix}
    \qquad
    A = \begin{bmatrix}
       2 & -3 &  0 &  1 \\
      -3 &  0 & -1 &  2 \\
       0 & -1 &  4 &  5 \\
       1 &  2 &  5 & -1
    \end{bmatrix}
    \qquad
    B = \begin{bmatrix} 2 \\ -1 \\ 0 \\ -3 \end{bmatrix}
    \qquad
    f = -10
  \]
\end{question}

%------------------------------------------------------------------------------%
\begin{resolution}{Solução}
  \[
    \onslide<+->{v^tAv + B^tv + f = 0}
  \]
  \[
    \onslide<+->{
    \begin{bmatrix} x & y & z & w \end{bmatrix}
    \begin{bmatrix}
       2 & -3 &  0 &  1 \\
      -3 &  0 & -1 &  2 \\
       0 & -1 &  4 &  5 \\
       1 &  2 &  5 & -1
    \end{bmatrix}
    \begin{bmatrix} x \\ y \\ z \\ w \end{bmatrix}
    +
    \begin{bmatrix} 2 \\ -1 \\ 0 \\ -3 \end{bmatrix}
    \begin{bmatrix} x \\ y \\ z \\ w \end{bmatrix}
    + (-10)
    = 0}
  \]
  \begin{gather*}
    \onslide<+->{2x^2 + 2(-3) xy + 2\times 0xz + 2\times 1xw}
    \onslide<+->{+ 0y^2 + 2(-1)yz + 2\times 2 yw} \\
    \onslide<+->{+ 4z^2 + 2\times 5zw + (-1)w^2}
    \onslide<+->{+ 2x + (-1)y + 0z + (-3)w}
    \onslide<+->{- 10}
    \onslide<+->{= 0}
  \end{gather*}
  \[
    \onslide<+->{2x^2 -6xy + 2xw}
    \onslide<+->{-2yz + 4yw}
    \onslide<+->{+ 4z^2 + 10zw -w^2}
    \onslide<+->{+ 2x -y -3w}
    \onslide<+->{= 10}
  \]
\end{resolution}

%------------------------------------------------------------------------------%
\endinput

%------------------------------------------------------------------------------%
\begin{question}
  Escreva a equação
  \[
    3x^2 + 2xy + 2y^2 - 4x + 6y = 1
  \]
  na forma matricial
  \(
    \vec{r}^{\,t}A\vec{x} + \vec{p}^{\,t}\vec{r} + f = 0
  \)

  \bigskip
  Determine $A$, $p$ e $f$
\end{question}

%------------------------------------------------------------------------------%
\begin{resolution}{Solução}
\begin{columns}[t]
\column{0.5\textwidth}
  Equação
  \[
    3x^2 + 2xy + 2y^2 - 4x + 6y = 1
  \]

  \pause
  Parte quadrática
  \[
    \pause
    ax^2 + 2bxy + cy^2
    \pause
    = 3x^2 + 2xy + 2y^2
  \]
  \[
    \pause
    A = \begin{bmatrix}
      3 & 1 \\
      1 & 2
    \end{bmatrix}
  \]

\column{0.4\textwidth}
  \pause
  Termos lineares
  \[
    \pause
    -4x + 6y
    \pause
    =
    \begin{bmatrix} -4 & 6  \end{bmatrix}
    \begin{bmatrix}  x \\ y \end{bmatrix}
  \]
  \[
    \pause
    p = \begin{bmatrix} -4 \\ 6 \end{bmatrix}
  \]

  \bigskip
  \pause
  Termo constante
  \[
    \pause
    f = -1
  \]
\end{columns}
\end{resolution}

%------------------------------------------------------------------------------%
\endinput


%------------------------------------------------------------------------------%
\begin{question}
  Escreva a equação
  \[
    2x - 3y + z = 5
  \]
  na forma matricial
  \(
    r^tAr + p^tr + f = 0
  \)
\end{question}

%------------------------------------------------------------------------------%
\begin{resolution}{Solução}
\begin{columns}[t]
\column{0.5\textwidth}
  Equação
  \[
    2x - 3y + z = 5
  \]

  \pause
  Parte quadrática
  \[
    A = \begin{bmatrix}
      0 & 0 \\
      0 & 0
    \end{bmatrix}
  \]

\column{0.4\textwidth}
  \pause
  Termos lineares
  \[
    \pause
    2x - 3y =
    \pause
    \begin{bmatrix} 2 & -3 \end{bmatrix}
    \begin{bmatrix} x \\ y \end{bmatrix}
  \]
  \[
    \pause
    p = \begin{bmatrix} 2 \\ -3 \end{bmatrix}
  \]

  \bigskip
  \pause
  Termo constante
  \[
    f = -5
  \]
\end{columns}
\end{resolution}

%------------------------------------------------------------------------------%
\endinput

%------------------------------------------------------------------------------%
\begin{question}
  Represente matricialmente a equação
  \[
    x^2 + y^2 = 9
  \]
\end{question}

%------------------------------------------------------------------------------%
\begin{resolution}{Solução}
\begin{columns}[t]
\column{0.5\textwidth}
  Equação
  \[
    x^2 + y^2 = 9
  \]

  \pause
  Parte quadrática
  \[
    \pause
    ax^2 + 2bxy + cy^2
    \pause
    = x^2 + 2y^2
  \]
  \[
    \pause
    A = \begin{bmatrix}
      1 & 0 \\
      0 & 2
    \end{bmatrix}
  \]

\column{0.4\textwidth}
  \pause
  Termos lineares
  \[
    \pause
    p = \begin{bmatrix} 0 \\ 0 \end{bmatrix}
  \]

  \bigskip
  \pause
  Termo constante
  \[
    f = -9
  \]
\end{columns}
\end{resolution}

%------------------------------------------------------------------------------%
\endinput

\begin{frame}{Exemplo 3 — Enunciado}
Represente matricialmente a equação
\[
x^2+y^2=9
\]
\end{frame}

\begin{frame}{Exemplo 3 — Resolução}
Definimos
\[
A=\begin{bmatrix}1&0\\0&1\end{bmatrix}
\qquad
\vec{x}=\begin{bmatrix}x\\y\end{bmatrix}
\]
Então
\[
\vec{x}^{\,t}A\vec{x}=9
\]
pois
\[
\begin{bmatrix}x&y\end{bmatrix}
\begin{bmatrix}1&0\\0&1\end{bmatrix}
\begin{bmatrix}x\\y\end{bmatrix}
=x^2+y^2
\]
A equação é a de uma circunferência centrada na origem.
\end{frame}


%------------------------------------------------------------------------------%
\section{Soluções de Equações de Segundo Grau em $\R^n$}

%------------------------------------------------------------------------------%
\begin{frame}{Equações degeneradas}
  \onslide<+->{Estamos interessados nas soluções das equações de segundo grau}

  \onslide<+->{Em \(\R^2\) elas são chamadas \emph{cônicas}}
  \begin{itemize}[<+->]
    \item elipses (circunferências)
    \item parábolas
    \item hipérboles
  \end{itemize}

  \onslide<+->{Em \(\R^3\) elas são chamadas \emph{quádricas}}
\end{frame}

%------------------------------------------------------------------------------%
\begin{frame}{Equações Degeneradas}
  Algumas equações de segundo grau representam curvas degeneradas

  \pause
  \bigskip
  Mesmo que a expressão envolva termos de segundo grau

  a solução é uma reta, ponto ou vazia
\end{frame}

%------------------------------------------------------------------------------%
%------------------------------------------------------------------------------%
\begin{question}
  Descreva o conjunto de soluções reais da equação
  \[
    x^2 + y^2 = 0
  \]
\end{question}

%------------------------------------------------------------------------------%
\begin{resolution}{Solução}
\begin{columns}
\column{0.4\textwidth}
  Como
  \[
    x^2\geq 0
    \qquad
    y^2\geq 0
  \]
  \pause
  a igualdade
  \[
    x^2 + y^2 = 0
  \]
  \pause
  só pode ocorrer quando
  \begin{align*}
    \onslide<+->{x^2 & = 0 & y^2 & = 0} \\
    \onslide<+->{x   & = 0 & y   & = 0}
  \end{align*}
  \column{0.4\textwidth}
  \onslide<+->{O conjunto de soluções é
  \[
    s = \big\{ (0,0) \big\}
  \]}
\end{columns}
\end{resolution}

%------------------------------------------------------------------------------%
\endinput

%------------------------------------------------------------------------------%
\begin{question}
  Descreva o conjunto de soluções reais da equação
  \[
    x^2 - y^2 = 0
  \]
\end{question}

%------------------------------------------------------------------------------%
\begin{resolution}{Solução}
\begin{align*}
  \onslide<+->{x^2 - y^2 & = 0  \\}
  \onslide<+->{(x-y)(x+y) & = 0  \\[4mm]}
  \onslide<+->{x - y       = 0}
  \onslide<+->{\qquad & \text{ou}\qquad  x + y  = 0 \\}
  \onslide<+->{y           = x}
  \onslide<+->{\qquad & \hphantom{\text{ou}\qquad  x + } y  = - x}
\end{align*}
\onslide<+->{Duas retas que passam pela origem}
\end{resolution}

%------------------------------------------------------------------------------%
\endinput

%------------------------------------------------------------------------------%
\begin{question}
  Descreva o conjunto de soluções reais da equação
  \[
  x^2+y^2+1=0
  \]
\end{question}

%------------------------------------------------------------------------------%
\begin{resolution}{Solução}
  Para todos os reais \(x\) e \(y\)
  \[
    x^2 + y^2 \geq 0
  \]
  \[
    \pause
    x^2 + y^2 + 1 \geq 1
  \]

  \pause
  A expressão não pode ser igual a zero

  \pause
  Portanto, a equação não possui soluções reais
\end{resolution}

%------------------------------------------------------------------------------%
\endinput


%%------------------------------------------------------------------------------%
%\section{Lista Mínima}
%
%%------------------------------------------------------------------------------%
%\begin{frame}{Lista Mínima}
%  \begin{enumerate}
%    \item
%  \end{enumerate}
%
%  \vfill
%  Atenção: A prova é baseada no livro, não nas apresentações
%\end{frame}

%------------------------------------------------------------------------------%
\end{document}
%------------------------------------------------------------------------------%
